% Part: history
% Chapter: biographies
% Section: alonzo-church

\documentclass[../../../include/open-logic-section]{subfiles}

\begin{document}

\olfileid[tr]{his}{bio}{chu}

\olsection{Alonzo Church}

\olphoto{church-alonzo}{Alonzo Church}

Alonzo Church, 14 Haziran 1903'te Washington, DC'de doğdu. Çocukluğunun
erken döneminde havalı tüfekle yaşanan bir kaza Church'ün bir gözünü kör
bıraktı. Connecticut'taki hazırlık okulunu 1920'de bitirdi ve aynı yıl
Princeton'da üniversite eğitimine başladı. Doktorasını 1927'de tamamladı.
Yurt dışında geçirdiği birkaç yılın ardından Princeton'a döndü. Church son
derece nazik ve dikkatli biri olarak bilinirdi. Tahtaya kusursuz yazar,
önemli belgeleri Duco çimentosuyla (saydam bir yapıştırıcıyla) özenle
kaplayarak korurdu. Akademik uğraşlarının dışında bilimkurgu dergileri
okumaktan hoşlanır, metinlerde yanlışlık görürse editörlere yazmaktan
çekinmezdi.

Church'ün akademik başarıları büyüktü. Öğrencileri Stephen Kleene ve
Barkley Rosser'la birlikte, Alan Turing'in Turing makinesini geliştirmesinden
bağımsız olarak etkin hesaplanabilirlik teorisini, lambda hesabını
geliştirdi. Hesaplanabilirliğin bu iki tanımı denktir ve günümüzde
\emph{Church--Turing Tezi} denen şu teze yol açar: doğal sayılar üzerindeki
bir fonksiyon, ancak ve ancak Turing makinesiyle (ya da lambda hesabıyla)
hesaplanabiliyorsa etkin biçimde hesaplanabilirdir. Church ayrıca günümüzde
\emph{Church Teoremi} denen sonucu ispatladı: birinci dereceden formüllerin
geçerliliğine ilişkin karar problemi çözülemez.

Church ileri yaşlarına kadar çalışmayı sürdürdü. 1967'de Princeton'dan
UCLA'ya geçti ve 1990'da emekli olana dek burada profesörlük yaptı. 1
Ağustos 1995'te 92 yaşında öldü.

\begin{reading}
Church'ün kısa bir yaşam öyküsü için \citet{EndertonND} kaynağına bakın.
Lambda hesabı ve Entscheidungsproblem (Church Tezi) hakkındaki özgün
yazıları \citet{Church1936,Church1936a} kaynaklarıdır.
\citet{Aspray1984}, Church'le 1930'lardaki Princeton matematik çevresi
hakkında yapılan bir söyleşiyi kaydeder. Church 1936'dan 1979'a kadar
\emph{Journal of Symbolic Logic} için bir dizi kitap değerlendirmesi yazdı.
Bunların tümü John MacFarlane'in internet sitesinde arşivlenmiştir
\citep{MacFarlane2015}.
\end{reading}
\end{document}
